\documentclass[
 reprint,
 amsmath,amssymb,
 aps,
 prl,
 superscriptaddress,
 nofootinbib
]{revtex4-2}

\usepackage{bm}
\usepackage{xcolor}
\usepackage[colorlinks=true,citecolor=blue!55!black,
linkcolor=blue!55!black,urlcolor=blue!55!black]{hyperref}

\newcommand{\Om}{\Omega}
\newcommand{\Res}{\operatorname*{Res}}
\newcommand{\ii}{\mathrm{i}}

\begin{document}

\title{Supplemental Material for ``Quasinormal-Mode Collider Spectroscopy''}

\author{Mian Zhu}
\email{zhumian@scu.edu.cn}
\affiliation{College of Physics, Sichuan University, Chengdu 610065, China}

\maketitle

\section{Perturbative origin of the quadratic QNM response}
\label{sec:quadratic-response}

Here we clarify the perturbative origin of the quadratic response used in the main text. Consider an expansion of the metric about a stationary black-hole background,
\begin{equation}
	g_{\mu\nu}
	=
	\bar g_{\mu\nu}
	+\epsilon h_{\mu\nu}^{(1)}
	+\epsilon^2 h_{\mu\nu}^{(2)}
	+\mathcal O(\epsilon^3).
\end{equation}
Let \(\mathcal E[g]=0\) denote the exact gravitational field equations. At first and second order, their schematic expansion is
\begin{align}
	\mathcal E^{(1)}[h^{(1)}] &= J, \\
	\mathcal E^{(1)}[h^{(2)}]
	&=
	S^{(2)}[h^{(1)},h^{(1)}],
	\label{eq:second-order-field-equation}
\end{align}
where \(J\) is a bookkeeping source for the linear response and
\(S^{(2)}\) collects terms quadratic in \(h^{(1)}\). Overall signs and numerical symmetry factors are included in the definition of \(S^{(2)}\).

Equation~\eqref{eq:second-order-field-equation} shows that the same linearized operator governs the propagation of \(h^{(1)}\) and \(h^{(2)}\). The nonlinearity is contained in the effective source \(S^{(2)}\), rather than in the propagation operator acting on \(h^{(2)}\). Thus, the appearance of the linear retarded Green function in the second-order solution is a direct consequence of perturbation theory and is not a projection of the nonlinear perturbation onto a linear solution.

We decompose \(h^{(2)}\) and \(S^{(2)}\) into the same asymptotic angular and parity channels used to define the observable strain multipoles. For a channel \(a\), the interaction-generated solution satisfying retarded boundary conditions is
\begin{equation}
	h_{a,\mathrm{dr}}^{(2)}(\omega)
	=
	G_a^R(\omega)S_a^{(2)}(\omega).
	\label{eq:retarded-second-order-solution}
\end{equation}
The general solution may additionally contain a homogeneous contribution,
\begin{equation}
	h_a^{(2)}
	=
	h_{a,\mathrm{hom}}^{(2)}
	+h_{a,\mathrm{dr}}^{(2)}.
\end{equation}
The homogeneous part satisfies the linearized source-free equation and represents an independently excited ordinary QNM contribution. It may be included in the fitted linear-mode amplitudes. In QCS we focus on \(h_{a,\mathrm{dr}}^{(2)}\), whose frequency and amplitude are fixed by the parent modes.

For reference, the linear response in channel \(a\) is
\begin{equation}
	h_a^{(1)}(\omega)
	=
	G_a^R(\omega)J_a(\omega).
\end{equation}
A QNM \(A=(a,n)\) is a pole of the channel-projected retarded Green function,
\begin{equation}
	G_a^R(\omega)
	=
	\frac{Z_A}{\omega-\Omega_A}
	+G_{a,\mathrm{reg}}^R(\omega),
	\qquad
	Z_A
	\equiv
	\underset{\omega=\Omega_A}{\operatorname{Res}}
	G_a^R(\omega),
	\label{eq:linear-qnm-pole}
\end{equation}
where \(\Omega_A\) is the complex QNM frequency, with
\(\operatorname{Im}\Omega_A<0\), and \(Z_A\) is the corresponding Green-function residue in the adopted asymptotic normalization.

The channel-\(a\) component of the quadratic source can be written as
\begin{align}
	S_a^{(2)}(\omega)
	={}&
	\sum_{bc}
	\int\frac{d\omega_1}{2\pi}
	\frac{d\omega_2}{2\pi}\,
	2\pi\delta(\omega-\omega_1-\omega_2)
	\nonumber\\
	&\times
	\Gamma^{(2)}_{a|bc}(\omega;\omega_1,\omega_2)
	h_b^{(1)}(\omega_1)h_c^{(1)}(\omega_2).
	\label{eq:quadratic-source-kernel}
\end{align}
Here \(\Gamma^{(2)}_{a|bc}\) is the quadratic source kernel that couples the parent channels \(b,c\) to the output channel \(a\). It contains the angular projections, radial overlap, derivatives, and reconstruction terms appropriate to the chosen perturbation variables. We absorb into \(\Gamma^{(2)}_{a|bc}\) any symmetry factor associated with \(b\leftrightarrow c\). The delta function follows from time-translation invariance of the stationary background.

Combining Eqs.~\eqref{eq:retarded-second-order-solution} and
\eqref{eq:quadratic-source-kernel} gives the quadratic response quoted in the main text,
\begin{align}
	h_a^{(2)}(\omega)
	={}&
	G_a^R(\omega)
	\sum_{bc}
	\int\frac{d\omega_1}{2\pi}
	\frac{d\omega_2}{2\pi}\,
	2\pi\delta(\omega-\omega_1-\omega_2)
	\nonumber\\
	&\times
	\Gamma^{(2)}_{a|bc}(\omega;\omega_1,\omega_2)
	h_b^{(1)}(\omega_1)h_c^{(1)}(\omega_2).
	\label{eq:quadratic-response}
\end{align}

To expose its pole structure, place the parent modes
\(B=(b,n_B)\) and \(C=(c,n_C)\) on their QNM poles. In the time domain,
\begin{equation}
	h_B^{(1)}(t)
	=
	A_B e^{-i\Omega_B(t-t_*)},
	\qquad
	h_C^{(1)}(t)
	=
	A_C e^{-i\Omega_C(t-t_*)},
	\qquad t>t_*.
\end{equation}
Their quadratic product therefore sources the composite complex frequency
\begin{equation}
	\Omega_Q\equiv\Omega_B+\Omega_C.
\end{equation}
Assuming that the quadratic kernel is analytic near the parent poles, define its on-shell value by
\begin{equation}
	\Gamma^{(2)}_{a|BC}
	\equiv
	\Gamma^{(2)}_{a|bc}
	\bigl(\Omega_Q;\Omega_B,\Omega_C\bigr).
\end{equation}
The causal Fourier transform of the quadratic source then contains a simple pole at \(\omega=\Omega_Q\), and the interaction-generated response takes the form
\begin{equation}
	h_a^{(2)}(\omega)
	\simeq
	\frac{
		G_a^R(\omega)\Gamma^{(2)}_{a|BC}A_BA_C
	}{
		\omega-\Omega_Q
	}
	+\text{terms regular at }\omega=\Omega_Q.
	\label{eq:composite-pole}
\end{equation}
A frequency-independent phase associated with the Fourier-transform convention is absorbed into the definition of \(\Gamma^{(2)}_{a|BC}\). Taking the residue of Eq.~\eqref{eq:composite-pole} gives
\begin{equation}
	\underset{\omega=\Omega_B+\Omega_C}{\operatorname{Res}}
	h_a^{(2)}(\omega)
	=
	G_a^R(\Omega_B+\Omega_C)
	\Gamma^{(2)}_{a|BC}A_BA_C.
	\label{eq:quadratic-residue}
\end{equation}

Equation~\eqref{eq:quadratic-residue} separates three physically distinct ingredients: the excitation \(A_BA_C\) of the parent modes, the quadratic coupling \(\Gamma^{(2)}_{a|BC}\), and the propagation of the generated source through the output Green function \(G_a^R\). If \(\Omega_B+\Omega_C\) approaches a QNM frequency of the output channel, the latter produces an additional resonant enhancement. At exact degeneracy, the product of the source pole and the output pole must instead be interpreted through its causal time-domain limit, which yields the familiar secular response proportional to
\(t e^{-i\Omega_Q t}\).

\section{Retarded-response conventions}

We first fix the conventions needed to compare nonlinear response
coefficients. We use
\begin{equation}
X(t)=\int_{\mathcal C}\frac{d\omega}{2\pi}\,
e^{-\ii\omega t}X(\omega),
\label{eqS:fourier}
\end{equation}
with the retarded contour chosen for decaying quasinormal modes (QNMs),
\(\operatorname{Im}\Om_A<0\). Lower-case labels \(a,b,\ldots\) denote
asymptotic angular and parity channels. Capital labels \(A=(a,n)\) identify
individual QNM poles, with \(n\) the overtone number.

The external observable is the source-frame multipolar strain at future null
infinity,
\begin{equation}
h_+-\ii h_\times
=\frac{M_z}{D_L}
\sum_{\ell mn} A_{\ell mn}
e^{-\ii\Om_{\ell mn}(t-t_\star)}
{}_{-2}Y_{\ell m}.
\label{eqS:strain}
\end{equation}
Other perturbation variables may be convenient internally, but their response
coefficients must be translated to this fixed observable before comparison.

\section{Composite pole at second order}

To display the pole structure without gauge-dependent details, consider the
reduced causal system
\begin{equation}
D_a(\ii\partial_t)q_a(t)
=J_a(t)+\frac12\sum_{bc}
\lambda_{a|bc}q_b(t)q_c(t),
\label{eqS:model}
\end{equation}
where \(\lambda_{a|bc}=\lambda_{a|cb}\). The linear retarded propagator is
\(G_a^R(\omega)=D_a^{-1}(\omega)\), and near a QNM pole \(A=(a,n)\),
\begin{equation}
G_a^R(\omega)
=\frac{Z_A}{\omega-\Om_A}+G_{a,\mathrm{reg}}^R(\omega).
\label{eqS:pole}
\end{equation}
The first two perturbative solutions are
\begin{align}
q_a^{(1)}(\omega)
&=G_a^R(\omega)J_a(\omega),\\
q_a^{(2)}(\omega)
&=\frac12G_a^R(\omega)\sum_{bc}\lambda_{a|bc}
\int\frac{d\omega_1}{2\pi}\,
q_b^{(1)}(\omega_1)q_c^{(1)}(\omega-\omega_1).
\label{eqS:q2}
\end{align}
Differentiating twice with respect to the sources gives
\begin{equation}
\mathcal R^{(2)}_{a|bc}
=G_a^R(\omega)\,
\lambda_{a|bc}\,
G_b^R(\omega_1)G_c^R(\omega_2),
\label{eqS:R2}
\end{equation}
where \(2\pi\delta(\omega-\omega_1-\omega_2)\) is suppressed. Amputation of
the three linear responses recovers \(\lambda_{a|bc}\). In the full
gravitational problem, we denote the corresponding dressed amputated vertex by
\(\Gamma^{(2)}_{a|bc}\).

Place the two parent propagators near poles \(B\) and \(C\). Their convolution
then contains
\begin{equation}
\int\frac{d\omega_1}{2\pi}
\frac{Z_BZ_C}
{(\omega_1-\Om_B)(\omega-\omega_1-\Om_C)}
\propto
\frac{Z_BZ_C}{\omega-\Om_B-\Om_C}.
\label{eqS:convolution}
\end{equation}
The second-order signal therefore has a composite pole at
\(\Om_Q=\Om_B+\Om_C\). Absorbing the external residues into the on-shell
parent amplitudes gives
\begin{equation}
\Res_{\omega=\Om_B+\Om_C}h_a^{(2)}(\omega)
=G_a^R(\Om_B+\Om_C)\,
\Gamma^{(2)}_{a|BC}A_BA_C.
\label{eqS:factorization}
\end{equation}
The output factor is normally evaluated away from its own poles, so the
quadratic pole is composite rather than a third linear QNM pole. If
\(\Om_B+\Om_C\simeq\Om_A\), the output Green function supplies a separate
resonant enhancement.

\section{Mode normalization and quadratic benchmark}

Normalization introduces a second potential ambiguity. An internal modal
expansion \(h=\sum_A A_Au_A\) is unchanged by
\(u_A\to c_Au_A\), \(A_A\to A_A/c_A\). A bare overlap written using the
\(u_A\)'s is therefore normalization dependent. This ambiguity is removed
either by normalizing every mode to a specified outgoing strain coefficient
or by working with physical Green functions and amputating their external
responses. In either approach, the physical field and its source convention
must be fixed before the kernel is meaningful.

At second order, the conversion from a regularized
Regge--Wheeler--Zerilli master variable to transverse-traceless strain also
contains analytic quadratic reconstruction terms. Consequently, the
scheme-independent result is the complete dressed strain response
\begin{equation}
A_{440}^{(2)}
=\frac{0.919-0.0627\ii}{6}A_{220}^{2}
=0.15352e^{-0.06812\ii}A_{220}^{2}.
\label{eqS:benchmark}
\end{equation}
The factor \(1/6\) combines
\(
\left(\begin{smallmatrix}4&2&2\\-4&2&2\end{smallmatrix}\right)=1/3
\)
with the identical-parent factor \(1/2\). Equation~(\ref{eqS:benchmark})
reproduces the published Schwarzschild result~\cite{BucciottiEtAl2024PRD,
BucciottiEtAl2024JHEP}.

\section{Dressed exchange and its time-domain limit}

An exchanged damped mode is a specific QNM pole \(S=(s,n_S)\) that must be
produced at one quadratic vertex and converted at another. After amputation of
the external responses, the intermediate propagator remains:
\begin{equation}
\Gamma^{(3),\mathrm{ex}}_{a|bcd}
=\Gamma^{(2)}_{a|sd}\,G_s^R(\omega_b+\omega_c)\,
\Gamma^{(2)}_{s|bc}+\text{crossed channels}.
\label{eqS:exchange-propagator}
\end{equation}
Near the pole \(S\), using
\(G_s^R(\nu)\simeq Z_S/(\nu-\Om_S)\), the full amputated response becomes
\begin{equation}
\Gamma^{(3)}_{a|bcd}
=\Gamma^{(3),\mathrm{reg}}_{a|bcd}
+\frac{\Gamma^{(2)}_{a|Sd}Z_S\Gamma^{(2)}_{S|bc}}
{\omega_b+\omega_c-\Om_S}
+\text{crossed channels}.
\label{eqS:exchange}
\end{equation}
The uppercase \(S\) on a vertex leg denotes evaluation at the intermediate
pole; in the reduced constant-coupling model of Eq.~(\ref{eqS:model}), each
\(\Gamma^{(2)}\) reduces to the corresponding quadratic coupling \(\lambda\).
The pair-frequency pole is distinct from the total composite pole at
\(\omega=\omega_b+\omega_c+\omega_d\), the external parent poles, and a
possible output pole. Its residue factorizes:
\begin{equation}
\Res_{\omega_b+\omega_c=\Om_S}
\Gamma^{(3)}_{a|bcd}
=\Gamma^{(2)}_{a|Sd}Z_S\Gamma^{(2)}_{S|bc}.
\label{eqS:exchange-res}
\end{equation}

Restoring the output response and putting three parents \(B,C,D\) on shell,
let \(A=(a,n_A)\) denote a candidate output pole and define
\(\Om_T=\Om_B+\Om_C+\Om_D\),
\(\Delta_S=\Om_B+\Om_C-\Om_S\), and \(\Delta_A=\Om_T-\Om_A\). The exchange
contribution to the output amplitude is
\begin{equation}
A_a^{(3),\mathrm{ex}}
=G_a^R(\Om_T)\,
\frac{\Gamma^{(2)}_{a|SD}Z_S\Gamma^{(2)}_{S|BC}}{\Delta_S}
A_BA_CA_D.
\label{eqS:dressed-exchange}
\end{equation}
Near the output pole \(A\), this becomes
\begin{equation}
A_a^{(3),\mathrm{ex}}
\simeq
\frac{Z_AZ_S\Gamma^{(2)}_{A|SD}\Gamma^{(2)}_{S|BC}}
{\Delta_A\Delta_S}A_BA_CA_D.
\label{eqS:double-pole}
\end{equation}
This equation cleanly separates double-resonant kinematics from the dynamical
residue and overlap factors. It applies when the poles are isolated and the
vertices vary slowly over the detuning scale.

The apparent divergence in Eq.~(\ref{eqS:double-pole}) is regulated by the
causal initial-value problem. To see this explicitly, let
\(\Om_p=\Om_B+\Om_C\), \(\Om_u=\Om_S+\Om_D\), and define
\begin{equation}
F(\Om,\Om_A;t)
\equiv
\frac{e^{-\ii\Om t}-e^{-\ii\Om_A t}}{\Om-\Om_A}.
\label{eqS:F}
\end{equation}
For vanishing nonlinear initial data and frequency-independent vertices, the
sequentially forced output is proportional to
\begin{equation}
q_a^{(3),\mathrm{ex}}(t)
\propto
\frac{F(\Om_T,\Om_A;t)-F(\Om_u,\Om_A;t)}
{\Om_p-\Om_S}.
\label{eqS:finite-cascade}
\end{equation}
Away from degeneracy, Eq.~(\ref{eqS:finite-cascade}) reproduces the two
denominators in Eq.~(\ref{eqS:double-pole}). If only the internal frequencies
coalesce, the intermediate response contains
\begin{equation}
\frac{e^{-\ii\Om_p t}-e^{-\ii\Om_S t}}
{\Om_p-\Om_S}
\longrightarrow
-\ii t\,e^{-\ii\Om_S t}.
\label{eqS:single-secular}
\end{equation}
At exact double degeneracy,
\(\Om_p\to\Om_S\) and \(\Om_T\to\Om_A\), expansion of
Eq.~(\ref{eqS:finite-cascade}) gives
\begin{equation}
q_a^{(3),\mathrm{ex}}(t)
\propto
\frac{t^2}{2}e^{-\ii\Om_A t},
\label{eqS:double-secular}
\end{equation}
up to the overall phase, residues, and vertex factors. The quadratic secular
envelope is the time-domain image of two sequential simple poles. It does not
grow indefinitely: perturbation theory requires the dressed cubic correction
to remain smaller than the lower-order signal, while QNM damping and the
finite ringdown duration bound the observable interval.

\section{Kerr spectrum scan and branch continuation}

We search for pure-gravity exchange channels using Kerr modes with spin weight
\(s=-2\) and convention
\begin{equation}
M\Om_{\ell mn}=M\omega_{\ell mn}-\ii M\gamma_{\ell mn}.
\end{equation}
The exploratory scan takes fundamental parents with
\(2\leq\ell\leq4\), all allowed \(m\), intermediate modes
\(\ell_S\leq7\), and \(n_S\leq2\). We impose azimuthal conservation and the
spherical-limit triangle rules at both vertices. The leading near-extremal
chain is
\begin{equation}
210+220\rightarrow431,\qquad431+220\rightarrow651.
\label{eqS:chain}
\end{equation}

The public \texttt{qnm} sequences~\cite{Stein2019} terminate at
\(\chi=0.99\). To avoid extrapolating their interpolants, we start from the
cached roots at \(\chi=0.97\) and continue every branch to \(\chi=0.995\) in
steps \(\delta\chi=10^{-4}\). At each step, the previous frequency and
angular separation constant seed a fresh Leaver continued-fraction solve.
We use angular truncation \(\ell_{\max}=20\), at least \(600\) radial
continued-fraction terms, and polish the frequency to relative tolerance
\(10^{-11}\). The maximum estimated continued-fraction error across all
plotted roots is \(1.4\times10^{-15}\).

At \(\chi=0.99\), the independently polished modes are
\begin{align}
M\Om_{210}&=0.572748611-0.046208379\ii,\\
M\Om_{220}&=0.870892659-0.029390424\ii,\\
M\Om_{431}&=1.430709787-0.098718565\ii,\\
M\Om_{651}&=2.312989693-0.092483745\ii.
\label{eqS:frequencies}
\end{align}
They give
\begin{align}
M\Delta_S
&=M(\Om_{210}+\Om_{220}-\Om_{431})\nonumber\\
&=0.012931482+0.023119762\ii,\\
M\Delta_A
&=M(\Om_{210}+2\Om_{220}-\Om_{651})\nonumber\\
&=0.001544236-0.012505482\ii.
\label{eqS:detunings}
\end{align}
The internal-detuning minimum occurs at
\(\chi=0.9921\), where \(|M\Delta_S|=0.0257684\).
The output-detuning minimum occurs at
\(\chi=0.9884\), where \(|M\Delta_A|=0.0117775\).
The product \(|M\Delta_S\Delta_A|\) is minimized near \(\chi=0.9890\).

At the production and conversion vertices, respectively,
\begin{equation}
\begin{pmatrix}4&2&2\\-3&1&2\end{pmatrix}
=-0.235702,\qquad
\begin{pmatrix}6&4&2\\-5&3&2\end{pmatrix}
=-0.226455.
\end{equation}
Both are nonzero. Reflection parity is also allowed:
odd \(210\) plus even \(220\) produces odd \(431\), and odd \(431\) plus
even \(220\) produces odd \(651\). Kerr spheroidal mixing modifies the
numerical couplings without removing the azimuthal rule.

The scan establishes only the resonant kinematics. A physical enhancement also
requires the QNM residues \(Z_{431}\) and \(Z_{651}\), the two quadratic
gravitational vertices, and sufficient event-dependent excitation of the
\(210\) parent. Nearly extremal QNM branching~\cite{YangEtAl2013} makes
continued branch tracking essential when these dynamical quantities are
calculated.

\bibliography{ref}

\end{document}
